% !TeX program = xelatex
% Lesson 15: UTF-8 Chinese lecture; compile with latexmk -xelatex.
\documentclass[UTF8,11pt,a4paper]{ctexart}
\usepackage[left=21mm,right=21mm,top=20mm,bottom=22mm]{geometry}
\usepackage{amsmath,amssymb,hyperref,fvextra}
\setCJKmainfont{Noto Serif CJK SC}
\setCJKsansfont{Noto Sans CJK SC}
\setlength{\parindent}{2em}
\setlength{\parskip}{0.4em}
\hypersetup{colorlinks=true,urlcolor=blue,linkcolor=blue}
\fvset{breaklines=true,breakanywhere=true}
\begin{document}
\title{Lesson 15：首次调用 OpenAI Responses API}
\author{LLM Learn 大模型开发课程}
\date{2026 年 11 月}
\maketitle
\tableofcontents
\newpage
\section{教学目标与先修}
掌握安全密钥注入、SDK 请求与响应、Token 输出限制，并完成离线 Mock 验证。\par
先修：Lesson 14 HTTP 与 JSON。 预计 90–120 分钟。\par
\section{API 与模型计算的边界}
Python SDK 封装 HTTPS、鉴权与响应解析；远程服务执行 LLM 推理。调用 API 并不意味着训练或本地运行 Transformer 权重。

\section{Responses API 参数}
client.responses.create(model=...,input=...,max\_output\_tokens=...)；model 指目标模型，input 是输入内容，max\_output\_tokens 控制输出上限，output\_text 是可展示文本。

\section{保护密钥与费用}
OPENAI\_API\_KEY 仅由环境变量读取，不能写进源码或 Git 仓库。ChatGPT 订阅与 API 计费不同，建议先设预算。默认模型可用 OPENAI\_MODEL 覆盖，应按本人账户权限选择。

\section{验证策略}
测试使用 FakeResponses 记录 create 的参数，并返回固定的模拟 output\_text；这证明程序参数正确，不证明在线模型可用。真实 API 运行需使用本人有额度的 Key，记录真实的成功/失败日志。

\section{核心关系或公式}
\[
C\simeq\frac{N_{\mathrm{in}}P_{\mathrm{in}}+N_{\mathrm{out}}P_{\mathrm{out}}}{10^{6}}
\]
\section{本课完整源码：逐行中文注释}
下列程序与本课 code 目录的实际源文件完全对应。每个模块、函数及关键可执行语句都配有中文说明。
\begin{Verbatim}[fontsize=\scriptsize,breaklines=true,breakanywhere=true]
"""Lesson 15：安全地执行一次 OpenAI Responses API 调用，支持离线注入假客户端。"""

# 导入 os，通过环境变量管理 API Key 和模型名。
import os
# 导入 SDK 客户端和错误基类；仅在真实调用时访问服务端。
from openai import OpenAI, OpenAIError

# 教学默认模型可通过 OPENAI_MODEL 环境变量覆盖。
DEFAULT_MODEL = os.getenv("OPENAI_MODEL", "gpt-4.1-mini")


def ask_once(prompt: str, client=None, model: str = DEFAULT_MODEL) -> str:
    """向 Responses API 发送一次文本请求；可注入 mock 客户端；返回回复文本。"""
    # 拒绝空白输入，避免无意义的付费调用。
    if not prompt.strip():
        # 明确告知调用方输入不符合要求。
        raise ValueError("问题文本不能为空")
    # 只有真实客户端未注入时才需要读取 API Key。
    if client is None:
        # 环境变量不存在说明尚不具备在线调用条件。
        if not os.getenv("OPENAI_API_KEY"):
            # 不显示密钥，仅提示安全配置步骤。
            raise RuntimeError("请先配置 OPENAI_API_KEY 环境变量")
        # 创建 SDK 客户端，控制超时和网络失败时的重试次数。
        client = OpenAI(timeout=45, max_retries=1)
    # 向指定模型提交文本，限制最大输出 Token 以减少误用成本。
    response = client.responses.create(model=model, input=prompt, max_output_tokens=400)
    # SDK 的 output_text 是由响应中可显示文本提取的内容。
    return response.output_text


def main() -> None:
    """在命令行执行首次真实 API 调用，并处理配置及网络错误。"""
    # 本课的问题专门围绕下一阶段要学的 Transformer 概念。
    question = "用简明中文解释 Transformer 中自注意力的作用。"
    try:
        # 调用封装函数并获得模型文本回答。
        answer = ask_once(question)
    except (OpenAIError, RuntimeError, ValueError) as error:
        # 将配置或 SDK 异常说明打印出来，不泄露任何令牌。
        print("调用失败:", error)
        # 不把失败误写为在线成功验证。
        return
    # 成功时显示真实返回的文本内容。
    print("模型回答:", answer)


if __name__ == "__main__":
    # 只有直接运行该文件才发起可能收费的请求。
    main()
\end{Verbatim}
\section{练习与验收}
\begin{itemize}
\item 空问题在请求前被拒绝。
\item Mock 客户端收到正确的 model、input、max\_output\_tokens。
\item 不会提交 API Key；真实在线调用尚未验证时必须如实注明。
\end{itemize}
\section{执行记录与参考}
使用 python -m unittest discover -s tests -v 执行测试。请以本课 results/verification.md 中的真实结果和 GitHub Actions 日志为准，Mock 并非在线 API 验证。
\begin{itemize}
\item \url{https://git-scm.com/docs}
\item \url{https://docs.github.com/en/actions}
\item \url{https://developer.mozilla.org/en-US/docs/Web/HTTP}
\item \url{https://platform.openai.com/docs/api-reference/responses}
\end{itemize}
\end{document}
